% Generated draft: requires mathematical and source review.
\documentclass[11pt]{article}
\usepackage[margin=25mm]{geometry}
\usepackage{fontspec}
\setmainfont{Latin Modern Roman}
\usepackage{amsmath,amssymb,mathtools,unicode-math}
\setmathfont{Latin Modern Math}
\usepackage{enumitem,needspace,xurl,hyperref}
\setlength{\parindent}{0pt}
\setlength{\parskip}{5pt}
\setlength{\emergencystretch}{3em}
\newcommand{\N}{\mathbb{N}}
\newcommand{\Z}{\mathbb{Z}}
\newcommand{\Q}{\mathbb{Q}}
\newcommand{\R}{\mathbb{R}}
\newcommand{\C}{\mathbb{C}}
\newcommand{\F}{\mathbb{F}}
\newcommand{\A}{\mathbb{A}}
\newcommand{\PP}{\mathbb{P}}
\newcommand{\E}{\mathbb{E}}
\newcommand{\eps}{\varepsilon}
\newcommand{\dd}{\,\mathrm{d}}
\begin{document}
\section*{UnsolvedMath ID: 2301040}
\textbf{Source problem code:} catalogue entry
\subsection{Definitions and mathematical statement}
Let $\mathbb D=\{z\in\C:|z|<1\}$ and let $f$ be meromorphic in $\mathbb D$ with finite order. With the standard disk Nevanlinna notation, $T(r,f)$ is the characteristic, $N(r,a;f)$ is the integrated counting function of $a$-points with multiplicity, and $\overline N(r,a;f)$ counts distinct $a$-points. The usual disk order is
\[
 \rho=\limsup_{r\uparrow1}\frac{\log^+T(r,f)}{\log(1/(1-r))}.
\]
For functions with the required nondegenerate characteristic growth, define
\[
 \delta(a,g)=1-\limsup_{r\uparrow1}\frac{N(r,a;g)}{T(r,g)},\qquad
 k(f)=\limsup_{r\uparrow1}\frac{\overline N(r,\infty;f)}{T(r,f)+1}.
\]
The sum over finite deficient values means the supremum of finite partial sums over $a\in\C$.

The source attributes to Shea and Sons the inequality
\[
 \sum_{a\in\C}\delta(a,f)
 \leq \delta(0,f')\bigl(1+k(f)\bigr)+\frac{2}{\lambda}(\rho+1).
\]
It prints
\[
 \lambda(f)=\liminf_{r\to\infty}
 \frac{T(r,f)}{\log(1/(1-r))}.
\]
That limit is outside the radial domain $0<r<1$ and does not define the stated disk quantity. A plausible intended correction is $r\uparrow1$; it is explicitly conditional here, not silently made. This normalized-characteristic quantity must also not be confused with the lower order $\liminf\log^+T/\log(1/(1-r))$.

After recovering the correct definition and admissible range of $\lambda$, L. R. Sons asks whether the coefficient $2$ can be replaced by $1$, namely whether
\[
 \sum_{a\in\C}\delta(a,f)
 \leq \delta(0,f')\bigl(1+k(f)\bigr)+\frac{1}{\lambda}(\rho+1).
\]
The source says this improvement would be best possible. It does not state the needed $\lambda>0$ assumptions or conventions for $\lambda=0,\infty$, or the growth conditions ensuring the deficiencies of $f$ and $f'$ are defined. These must be resolved before the displayed conditional target becomes a complete theorem.
\subsection{Short English statement}
Sharpen the reported disk deficiency inequality by replacing the coefficient two by one in the order-dependent term. The source says that improvement would be sharp. The printed definition of $\lambda$ has an impossible radial limit for a disk function, so the intended boundary limit and parameter assumptions require review. The supplied triage says the problem was open in Hayman's 2018 update and no recent resolution was located; this draft makes no independent current-status determination.
\subsection{Sources}
\begin{itemize}
\item UnsolvedMath permanent ID 2301040, source code catalogue entry, ``Research Problems in Function Theory --- Problem 1.40.'' Registry curation and metadata: Ulam AI and UnsolvedMath Contributors, CC BY 4.0; updated 2026-08-06.
\item Hayman, \emph{Research Problems in Function Theory} (2018 source update), Problem 1.40, attributed to L. R. Sons: \url{https://arxiv.org/abs/1809.07200}. Supplied source-TeX extraction accessed 2026-07-29. The metadata retains the 2018 update for status triage and marks status/release and rights review NEEDS\_REVIEW.
\item Shea and Sons are named in the supplied statement for the coefficient-two inequality, without a complete bibliographic reference. No title, date, theorem number or locator is inferred.
\end{itemize}
\end{document}
